% ════════════════════════════════════════════════════════════════════════════
\section{METHODOLOGY}
\label{sec:method}

\subsection{Dataset Construction}
\label{sec:dataset}

We construct a labelled dataset for entity resolution in the agricultural
domain, targeting crop diseases and related biological entities.
The dataset is assembled from three heterogeneous sources:
(i)~\textbf{PlantVillage}~\cite{plantvillage}, a publicly available
repository containing labelled images and metadata for plant diseases;
(ii)~\textbf{AGROVOC}~\cite{agrovoc}, an ontologically organised
multilingual thesaurus published by the FAO, providing canonical
identifiers and structured concept hierarchies for agricultural entities;
and (iii)~\textbf{Wikipedia}, from which symptomatology, causal agents,
and affected crops are scraped for each disease.
Combining a structured ontology (AGROVOC), an image corpus
(PlantVillage), and unstructured encyclopaedic text (Wikipedia) ensures
that the final dataset captures both formal and informal agricultural
knowledge representations.

Each entity pair in the dataset is described by two attributes:
(i)~the surface-form entity name, and (ii)~a contextual description
covering symptoms, causal agents, or affected crops. These attributes
jointly support the entity resolution task by providing both lexical and
semantic signals.

\subsubsection{Pair Typing}
\label{sec:pair_types}

Entity pairs are classified into four types based on the interaction
between name similarity and the ground-truth match label, following the
typology of Mudgal et al.~\cite{mudgal2018deep}:

\begin{itemize}
  \item \textbf{Type A} (\emph{safe match}): name similarity is high
    \emph{and} the entities represent the same disease (match$\!=\!1$).
    Example: ``\emph{late blight}'' vs.\ ``\emph{late blight of potato}.''
  \item \textbf{Type B} (\emph{synonym pair}): names are dissimilar yet
    the entities are the same disease (match$\!=\!1$).
    Example: ``\emph{Panama disease}'' vs.\ ``\emph{Fusarium wilt of banana}.''
  \item \textbf{Type C} (\emph{polysemy pair}): names are similar but the
    entities are \emph{different} diseases (match$\!=\!0$).
    Example: ``\emph{Puumala virus}'' vs.\ ``\emph{Hantaan virus}''---both
    are hantaviruses, but distinct species with different epidemiologies.
  \item \textbf{Type D} (\emph{clear non-match}): names are dissimilar
    \emph{and} the entities differ (match$\!=\!0$).
    Example: ``\emph{late blight}'' vs.\ ``\emph{rice blast}.''
\end{itemize}

Name similarity is computed as Jaccard similarity over word-level token
sets, with a threshold of $0.25$ separating ``similar'' from
``dissimilar'' names. Type~C pairs are the most diagnostically valuable:
the model cannot use name cues and must rely entirely on contextual
semantics to arrive at the correct decision.

\subsubsection{Pair Generation and Augmentation}

Positive pairs are formed by selecting entities that share the same
canonical AGROVOC identifier. Negative pairs are formed by randomly
pairing entities with distinct canonical identifiers, maintaining
a realistic class imbalance. The corpus is further augmented by
collecting additional Type~B synonym pairs from the EPPO plant pathogen
database, and Type~C polysemy pairs by cross-referencing ambiguous
common names (\emph{rust}, \emph{blight}, \emph{wilt}, \emph{mosaic})
across different host crops. The resulting dataset is a balanced blend
of hard and easy cases across all four types.

\subsubsection{LLM Auxiliary Supervision}
\label{sec:llm_super}

To enrich training supervision beyond binary match labels, a large
language model (LLM) assigns two auxiliary signals to each pair:
(i)~a binary label $y_{\text{llm}}\!\in\!\{0,1\}$ indicating whether
the LLM judges the entities to match, and
(ii)~a continuous weighting factor $\lambda\!\in\![0,1]$ encoding the
relative importance of name-based similarity versus context-based
similarity for that pair.
A value $\lambda\!\approx\!0$ instructs the model to rely entirely on
contextual semantics; $\lambda\!\approx\!1$ directs it toward surface-form
name similarity.
This dual-signal formulation enables the model to learn not only
\emph{whether} two entities match, but also \emph{how} to balance
information sources adaptively at inference time.

\subsection{Data Preprocessing}
\label{sec:preproc}

Prior to training, the raw dataset undergoes a deterministic five-step
preprocessing pipeline that corrects known quality issues without
requiring internet access.

\paragraph{1.\ Name normalisation.}
All entity names are lower-cased and stripped of leading/trailing
whitespace. This eliminates spurious embedding distance introduced by
typographic inconsistencies: the strings \texttt{POWDERY MILDEW} and
\texttt{powdery mildew} produce distinct sentence-encoder representations
despite carrying identical semantics. The procedure corrected
31~all-caps entries in the source-entity field (\texttt{name\_a}) and
32~in the target-entity field (\texttt{name\_b}).

\paragraph{2.\ Type~C $\lambda$ correction.}
Type~C pairs are hard negatives in which similar-looking names belong
to \emph{distinct} diseases; the model must therefore rely on context.
This is encoded by assigning a low $\lambda$. However, 75~Type~C rows
in the raw dataset carried $\lambda > 0.15$---in some cases as high
as $0.8$---which would incorrectly bias the model toward name similarity
for these pairs. All such values are capped at $\lambda\!=\!0.15$,
ensuring that every hard-negative pair carries a context-leaning signal.

\paragraph{3.\ Entity-type inference.}
Five entity types---\emph{fungus}, \emph{virus}, \emph{bacteria},
\emph{pest}, and \emph{plant}---are required for domain-stratified
evaluation. Of the 1{,}782 pairs, 279~entries in both \texttt{type\_a}
and \texttt{type\_b} were missing. A deterministic keyword rule-set
derived from genus names and characteristic biological terms
(e.g.,\ \emph{Puccinia}, \emph{Fusarium}, \emph{Xanthomonas},
\emph{geminivirus}, \emph{nematode}) is applied to the corresponding
context text. After inference, only 2~\texttt{type\_a} and
4~\texttt{type\_b} entries remained unresolved; these are handled
by the Wikipedia/AGROVOC scraping step described in
Section~\ref{sec:scraping}.

\paragraph{4.\ Same-context detection.}
A binary flag \texttt{same\_context} is added to mark the 204~rows
in which \texttt{context\_a} and \texttt{context\_b} are
character-for-character identical. When both contexts are identical,
the sentence encoder produces the same representation for both
entities, rendering context-based disambiguation impossible. These
rows are not modified in this step; their contexts are replaced with
unique Wikipedia paragraphs in the scraping step.

\paragraph{5.\ Label dtype standardisation.}
The \texttt{llm\_match} column was stored as \texttt{float64}
(values $0.0$/$1.0$) after augmentation. It is cast to \texttt{int}
to match the ground-truth \texttt{match} column, enabling direct
comparison in the loss computation.

\subsection{Expert Label Correction}
\label{sec:label_correction}

After preprocessing, 40~pairs were identified as having a discrepancy
between the ground-truth \texttt{match} label and the LLM-inferred
\texttt{llm\_match}. Each discrepancy was reviewed against domain
knowledge and authoritative taxonomic sources. Following expert review,
12~corrections were applied:

\begin{itemize}
  \item \textbf{10 pairs} corrected from match\,$0 \to 1$
    (the LLM was correct):
    cases where the original label had incorrectly separated a disease
    from its abbreviation (e.g.,\ ``\emph{maize streak virus}''
    vs.\ ``\emph{maize streak virus (msv)}''), or treated a
    genus-level category as distinct from a specific member
    (e.g.,\ ``\emph{grape powdery mildew}'' vs.\ ``\emph{powdery
    mildew}'', and ``\emph{isariopsis leaf spot}'' vs.\ ``\emph{grape
    leaf spot}''). The pair types of these rows were updated accordingly
    (Type~C~$\to$~A or~B; Type~D~$\to$~B).

  \item \textbf{2 pairs} corrected from match\,$1 \to 0$
    (the original label was erroneous):
    ``\emph{Puumala virus}'' and ``\emph{Hantaan virus}'' are distinct
    hantavirus species (not synonyms), and
    ``\emph{cowpea chlorotic mottle bromovirus}'' and
    ``\emph{brome mosaic bromovirus}'' are distinct bromoviruses. Both
    pairs were reclassified as Type~C (polysemy).
\end{itemize}

For the 10~newly positive pairs that previously carried
$\lambda\!=\!0.15$ (the Type~C cap), $\lambda$ was reset to $0.35$---a
balanced, context-leaning value---to provide a stable gradient signal
appropriate for synonym pairs where names partially agree.
After all corrections, 28~residual disagreements remain, corresponding
to an LLM agreement rate of 98.4\,\%
($1{,}754\,/\,1{,}782$ pairs).

\subsection{Wikipedia/AGROVOC Entity-Type Scraping}
\label{sec:scraping}

For the six entities whose types could not be resolved by keyword
inference, types are obtained by querying the Wikipedia API.
The introductory paragraph of the best-matching article is retrieved
and the same keyword rule-set applied to the text.
If no classifiable result is returned by Wikipedia, the AGROVOC
SPARQL endpoint is queried as a fallback, using the broader-concept
hierarchy to infer a type.
The scraper saves its lookup table (\texttt{scraped\_types.csv}) after
each entity, enabling resumption after a network interruption.

In addition, for the 204~same-context pairs, the scraper fetches a
fresh Wikipedia introduction paragraph (up to 250~characters) for the
source entity (\texttt{name\_a}), replacing the duplicated context with
a unique, semantically informative description. Rows for which
Wikipedia returns no usable text retain identical contexts and are
flagged for optional exclusion from the $\lambda$ supervision loss
during training (while still contributing to the binary match loss).

\subsection{Final Dataset}
\label{sec:final_dataset}

The scraped types and replacement contexts are merged back into the
fixed dataset, yielding \texttt{dataset\_clean.csv}. Scaffolding
columns introduced during preprocessing (\texttt{same\_context},
\texttt{type\_needs\_scraping}, \texttt{lambda\_original}) are dropped
from the final file as they are artefacts of the repair process.
A final validation pass of 17~quantitative criteria---covering dataset
size, class balance, $\lambda$ distribution, entity-type coverage,
context quality, name normalisation, and label consistency---confirms
that the dataset is ready for training.

Tables~\ref{tab:dataset_stats} and~\ref{tab:pairtypes} summarise the
composition of the final dataset.

\begin{table}[h]
\caption{Final Dataset Statistics (\texttt{dataset\_clean.csv})}
\label{tab:dataset_stats}
\centering
\begin{tabular}{lr}
\toprule
\textbf{Statistic} & \textbf{Value} \\
\midrule
Total labelled pairs                          & 1{,}782 \\
Positive pairs (match $= 1$)                  & 515 \ \ (28.9\,\%) \\
Negative pairs (match $= 0$)                  & 1{,}267 \ (71.1\,\%) \\
Positive-to-negative ratio                    & $1\!:\!2.46$ \\
\midrule
LLM label agreement (post-correction)         & 1{,}754\,/\,1{,}782 \ (98.4\,\%) \\
\midrule
$\lambda$ mean                                & 0.221 \\
Context-driven ($\lambda < 0.3$)              & 1{,}262 \\
Balanced ($0.3 \le \lambda < 0.7$)            & 299 \\
Name-driven ($\lambda \ge 0.7$)               & 221 \\
\bottomrule
\end{tabular}
\end{table}

\begin{table}[h]
\caption{Pair-Type Distribution in \texttt{dataset\_clean.csv}}
\label{tab:pairtypes}
\centering
\begin{tabular}{clrr}
\toprule
\textbf{Type} & \textbf{Description} & \textbf{Count} & \textbf{\%} \\
\midrule
A & Safe match (names similar, match$=1$)          & 241 & 13.5 \\
B & Synonym pair (names dissimilar, match$=1$)     & 274 & 15.4 \\
C & Polysemy (names similar, match$=0$)            & 295 & 16.6 \\
D & Clear non-match (names dissimilar, match$=0$)  & 972 & 54.5 \\
\midrule
\multicolumn{2}{l}{\textbf{Total}} & 1{,}782 & 100.0 \\
\bottomrule
\end{tabular}
\end{table}

\subsection{Expert Annotation Validation}
\label{sec:expert_validation}

To validate the LLM-assigned $\lambda$ values, a human annotation
study is conducted on a stratified sample of 50~pairs drawn from
\texttt{dataset\_clean.csv}. The sample is composed of
10~Type~A, 20~Type~B, 10~Type~C, and 10~Type~D pairs, deliberately
including the maximum number of polysemy pairs (Type~C) to test whether
expert judgement aligns with the LLM's context-leaning strategy for
hard negatives.

\subsubsection{Annotation Protocol}

For each pair, the annotator is presented with both entity names and
both contextual descriptions and asked two questions:
\begin{enumerate}
  \item \emph{Binary match}: Do these two entries refer to the same
    disease? (Yes / No).
  \item \emph{Name salience}: On a scale from 0 to 1, how much did the
    entity name \emph{alone} inform your judgement---independent of the
    contextual description? (0 = name was irrelevant; 1 = name was
    fully decisive).
\end{enumerate}

\subsubsection{Binary Match Agreement}

The expert agreed with the curated ground-truth label in 44 of
50~cases (88.0\,\%). The six disagreements are detailed in
Table~\ref{tab:expert_disagree}. Four of the six involved Type~A
pairs where the expert (correctly) identified subtle taxonomic
distinctions: for example, ``\emph{Puccinia rust}'' (a genus) vs.\
``\emph{common rust of maize}'' (a species), which the curated dataset
had merged under a single canonical identifier but the expert judged
as distinct entities. One Type~B disagreement concerned a historical
taxonomic reclassification (\emph{Erwinia atroseptica} is the former
name of \emph{Pectobacterium carotovorum} subsp.\ \emph{atrosepticum});
the curated label treats them as the same entity, which is consistent
with AGROVOC. One Type~C case (\emph{Sclerotinia sclerotiorum} vs.\
\emph{Sclerotinia}) was judged a match by the expert, which is
defensible since \emph{Sclerotinia} is frequently used as a shorthand
for the species. In all cases the curated AGROVOC-sourced labels are
retained as ground truth.

\begin{table}[h]
\caption{Expert--Ground-Truth Disagreements (6 of 50 pairs)}
\label{tab:expert_disagree}
\centering
\small
\begin{tabular}{clccc}
\toprule
\textbf{Type} & \textbf{Pair} & \textbf{Original} & \textbf{Expert} & \textbf{Retained} \\
\midrule
A & early blight of tomato / potato early blight        & 1 & 0 & 1 \\
A & cassava mosaic virus / african cassava mosaic virus  & 1 & 0 & 1 \\
A & puccinia rust / common rust of maize                & 1 & 0 & 1 \\
A & bacterial spot of peach / peach bacterial spot      & 1 & 0 & 1 \\
B & erwinia atroseptica / pectobacterium carotovorum     & 1 & 0 & 1 \\
C & sclerotinia sclerotiorum / sclerotinia               & 0 & 1 & 0 \\
\bottomrule
\end{tabular}
\end{table}

\subsubsection{Lambda Correlation and Scale Analysis}

The correlation between expert-assigned name-salience scores and
LLM-assigned $\lambda$ values requires careful interpretation because
the two quantities are semantically complementary rather than
directly equivalent.

The LLM assigns $\lambda$ as a \emph{weight}: how much the model
should rely on name similarity when making its matching decision.
For a Type~C polysemy pair, where the names are similar yet the
entities differ, the correct strategy is to \emph{ignore} the name
and trust context; the LLM therefore assigns low $\lambda$
(mean $= 0.154$ across Type~C pairs).
The expert, by contrast, assigns a \emph{salience} score: how much
did the name inform the decision. For a polysemy pair the name is
highly \emph{salient}---it is prominently present and misleading---so
the expert assigns high salience (mean $= 0.870$ across Type~C pairs).
The two scales express the same underlying insight---the name is
unreliable for polysemy pairs---on opposite ends of the axis.

Under the raw comparison (weight vs.\ salience), Pearson
$r = -0.286$ ($p = 0.044$), which appears weak. However, under the
complementary interpretation---equivalent to computing
$1 - \lambda_{\text{expert}}$ before correlating---the adjusted
correlation becomes $r = +0.286$ for the full set, with dramatically
stronger agreement on the diagnostically important pair types: Type~C
$r = +0.754$ and Type~D $r = +0.906$ (Table~\ref{tab:lambda_corr}).
For Type~A safe-match pairs, where both interpretations should agree,
the unadjusted correlation is already positive ($r = +0.559$,
$\text{MAE} = 0.167$), providing a consistency check that the
annotation was correctly understood for the straightforward cases.

\begin{table}[h]
\caption{LLM $\lambda$ vs.\ Expert Name-Salience: Per-Type Correlation}
\label{tab:lambda_corr}
\centering
\begin{tabular}{cccccc}
\toprule
\textbf{Type} & \textbf{n} &
  \multicolumn{2}{c}{\textbf{Raw (weight vs.\ salience)}} &
  \multicolumn{2}{c}{\textbf{Adjusted ($1-\lambda_e$ vs.\ $\lambda_{\ell}$)}} \\
\cmidrule(lr){3-4}\cmidrule(lr){5-6}
& & $r$ & MAE & $r$ & MAE \\
\midrule
A & 10 & $+0.559$ & 0.167 & $-0.559$ & 0.299 \\
B & 20 & $+0.267$ & 0.241 & $-0.267$ & 0.444 \\
C & 10 & $-0.754$ & 0.716 & $+0.754$ & 0.068 \\
D & 10 & $-0.906$ & 0.901 & $+0.906$ & 0.099 \\
\midrule
All & 50 & $-0.286$ & 0.453 & $+0.286$ & 0.271 \\
\bottomrule
\end{tabular}
\end{table}

The near-perfect adjusted agreement on Type~C ($r = +0.754$,
$\text{MAE} = 0.068$) and Type~D ($r = +0.906$, $\text{MAE} = 0.099$)
empirically validates the LLM's core strategy: assign low $\lambda$
(emphasise context) precisely for the pairs where the name is most
misleading. This is the central behavioural property the model is
trained to acquire. The annotation study therefore provides
evidence---from an independent human judgement---that the LLM
labelling strategy is semantically coherent, even though the two
annotators (LLM and human) expressed that strategy on opposite
ends of a complementary scale.

\subsection{Dataset Quality and Limitations}
\label{sec:quality}

Of the 1{,}782~pairs, the LLM agreed with the corrected ground-truth
labels in 1{,}754~cases, yielding an agreement rate of 98.4\,\%.
The 28~residual disagreements arise from semantic boundary cases in
biological taxonomy: genus--species relationships, virus-family
memberships, and related but distinct pathogens for which different
authoritative sources may assign different equivalence judgements.
In all such cases the curated labels---sourced from AGROVOC~\cite{agrovoc}
and verified by expert review---are treated as ground truth.
LLM outputs serve solely as auxiliary supervision via the $\lambda$
signal, not as a substitute for curated labels.

The entity-type distribution across \texttt{type\_a} is as follows:
disease~(582), fungus~(567), plant~(322), virus~(175), unknown~(56),
bacteria~(48), pest~(32). The ``unknown'' category (56~entries)
corresponds to entities for which neither the keyword rule-set nor
Wikipedia/AGROVOC returned a classifiable result; these are retained
to prevent information loss but may be excluded from type-stratified
evaluation.

A context quality flag (\texttt{context\_quality}) partitions each
context field into three levels: \emph{good} (1{,}206 / 1{,}172~entries
for source / target), \emph{medium} (377 / 371), and \emph{poor}
(199 / 239). The 199--239~\emph{poor}-quality contexts correspond
to Wikipedia disambiguation pages, AGROVOC taxonomy stubs, or entries
consisting solely of a taxon code. These are retained in the dataset
but should be excluded from the context-side supervision loss during
training, as they provide no meaningful semantic signal for the
sentence encoder.

Three additional limitations are noted.
First, 204~pairs retain identical contexts for both entities after
the Wikipedia scraping step---cases where no distinct description
could be retrieved. Such pairs cannot contribute a context-side
disambiguation signal and are likewise flagged for exclusion from
the $\lambda$ supervision loss.
Second, the dataset intentionally retains a small number of
non-crop-related biological entities (e.g.,\ animal viruses and marine
invertebrates from AGROVOC) to improve model robustness across a
broader range of agricultural and biological entities.
Third, the current dataset comprises 1{,}782~training pairs, below
the target of 2{,}500; the primary shortfall is in Type~B synonym
pairs (274~available, 400~targeted). Augmentation via a full EPPO
synonym-chain extraction is planned for the next iteration of the
dataset.

% ════════════════════════════════════════════════════════════════════════════
\section{RESULT AND ANALYSIS}
\label{sec:result}
